\section[Электростатическое кольцо: коэффициент скольжения и равновесный сдвиг]{Электростатическое кольцо:\\ коэффициент скольжения\\ и равновесный сдвиг}
\label{sec:es-base}

В электростатическом дефлекторе часть полной энергии частицы переходит в потенциальную, поэтому скорость внутри дефлектора зависит от радиального положения. Это меняет и дисперсию, и время пролёта. В главе предложена модель потенциала дефлектора, выведены коэффициенты запаздывания электростатического кольца и решены продольные уравнения. Бетатронные колебания не учитываются: частица движется по замкнутой орбите в плоскости кольца, $y=0$.

\subsection{Постановка и потенциал дефлектора}

Поворот пучка создают только электростатические дефлекторы; вне них $1/\rho=0$ и электрический потенциал равен нулю, а фокусирующие элементы магнитные. Резонатор расположен вне дефлекторов. Во всех дефлекторах один и тот же показатель поля $n$; в плоскости кольца
\begin{equation}
\label{eq:Ex}
E_x=E_0\Bigl(1-n\frac{x}{\rho}+O(x^2)\Bigr),
\end{equation}
где $E_0$ --- поле на референсной орбите, $x$ --- радиальное отклонение, $y$ --- вертикальное. В дефлекторе продольный масштабный множитель $h_s=1+x/\rho$.

Центростремительная сила на референсной орбите задаёт
\begin{equation}
\label{eq:E0}
qE_0=-\frac{p_sv_s}{\rho}.
\end{equation}
Потенциал $\Phi$ отсчитывается от референсной орбиты и измеряется в единицах $p_sv_s=\beta_s^2E_s$:
\begin{equation}
\label{eq:psi-def}
\chi=\frac{q\Phi}{p_sv_s}.
\end{equation}
Интегрирование $E_x=-\partial\Phi/\partial x$ с \eqref{eq:Ex}, \eqref{eq:E0} даёт члены по $x$, а уравнение Лапласа в координатах дефлектора (поле не зависит от $s$)
\begin{equation}
\label{eq:laplace}
\frac{\partial^2\chi}{\partial x^2}
+\frac{\partial^2\chi}{\partial y^2}
+\frac{1}{\rho+x}\frac{\partial\chi}{\partial x}=0
\end{equation}
фиксирует коэффициент при $y^2$:
\begin{equation}
\label{eq:psi}
\chi
=\frac{x}{\rho}
-\frac{n}{2}\frac{x^2}{\rho^2}
+\frac{n-1}{2}\frac{y^2}{\rho^2}
+O\bigl((x/\rho)^3,(y/\rho)^3\bigr).
\end{equation}
На внешнем радиусе $x>0$ потенциальная энергия частицы положительна. Член с $y^2$ на орбите в плоскости кольца не участвует; он нужен для вертикальных бетатронных колебаний.

\subsection{Полная и кинетическая энергия, скорость в дефлекторе}

Поле дефлектора статическое, поэтому между прохождениями резонатора сохраняется \emph{полная} энергия $\gamma mc^2+q\Phi$. \emph{Кинетическая} энергия $\gamma mc^2$ внутри дефлектора меняется вместе с потенциалом: при входе на внешний радиус часть её переходит в потенциальную, при выходе возвращается. Лоренц-фактор $\gamma$ всюду ниже --- кинетический.

Переменная $\delta$ определяется по \eqref{eq:delta-def} через импульс вне дефлекторов, где $\Phi=0$. Тем самым $\delta$ --- метка полной энергии, одинаковая вдоль всего оборота. Обозначим через $\gamma_{\mathrm{out}}(\delta)$ лоренц-фактор вне дефлекторов; он даётся рядом \eqref{eq:gamma-delta}. Сохранение полной энергии, $\gamma mc^2+q\Phi=\gamma_{\mathrm{out}}mc^2$, даёт локальный лоренц-фактор
\begin{equation}
\label{eq:gamma-loc}
\gamma=\gamma_{\mathrm{out}}(\delta)-\beta_s^2\gamma_s\chi
=\gamma_s+\beta_s^2\gamma_s(\delta-\chi)
+\frac{\beta_s^2}{2\gamma_s}\delta^2
+O(\delta^3).
\end{equation}
Член с $\chi$ точен, так как потенциал входит в полную энергию линейно.

Скорость в данной точке удобно записать через локальный импульс. Определим $\delta_{\mathrm{loc}}$ как значение $\delta$, при котором ряд \eqref{eq:gamma-delta} даёт локальный $\gamma$. Из \eqref{eq:gamma-loc}
\begin{equation}
\label{eq:delta-loc}
\delta_{\mathrm{loc}}
=\delta-\chi+\frac{2\delta\chi-\chi^2}{2\gamma_s^2}+O(3),
\end{equation}
где $O(3)$ --- члены суммарной третьей степени по $\delta$ и $\chi$. Это не новая динамическая переменная, а значение локального импульса в данной точке. Ряд \eqref{eq:vsv} в точке $\delta_{\mathrm{loc}}$ даёт
\begin{equation}
\label{eq:vsv-es}
\frac{v_s}{v}
=1-\frac{\delta-\chi}{\gamma_s^2}
+\frac{3(\delta-\chi)^2}{2\gamma_s^2}
-\frac{\delta^2}{2\gamma_s^4}
+O(3).
\end{equation}
На внешнем радиусе частица медленнее, чем в свободном промежутке при том же $\delta$.

Резонатор стоит при $\chi=0$, поэтому отображение \eqref{eq:map} и система \eqref{eq:smooth} сохраняют свой вид. От дефлекторов зависит только функция $\tau(\delta)$.

\subsection{Дисперсия в дефлекторе}
\label{sec:es-D}

Кривизну траектории задаёт произведение $pv$, то есть $1/R=-qE_x/(pv)$. Из \eqref{eq:gamma-loc} при $\chi=x/\rho$ и $pv=mc^2(\gamma-1/\gamma)$ в линейном порядке
\begin{equation}
\label{eq:cmp-pv}
\frac{pv}{p_sv_s}
=1+\Bigl(1+\frac{1}{\gamma_s^2}\Bigr)
\Bigl(\delta-\frac{x}{\rho}\Bigr).
\end{equation}
Подстановка в уравнение орбиты в сопровождающей системе, $x''+x/\rho^2=1/\rho-1/R$, даёт для замкнутой орбиты $x=D_1\delta$ внутри дефлектора
\begin{equation}
\label{eq:cmp-D-es}
D_1''+\frac{2-n+1/\gamma_s^2}{\rho^2}\,D_1
=\frac{1+1/\gamma_s^2}{\rho}.
\end{equation}
В жёсткости единица --- центробежный член, $-n$ --- градиент поля, а $1+1/\gamma_s^2$ --- уменьшение $pv$ на внешнем радиусе из-за потенциала: там то же поле поворачивает траекторию сильнее. Источник $(1+1/\gamma_s^2)/\rho$ вместо магнитного $1/\rho$ отражает то, что кривизну задаёт $pv$, а не $p$~\cite{Mane2008,Mane2012}. Вне дефлекторов уравнение переходит в $D_1''+K_1D_1=0$. На границах дефлектора $D_1$ и $D_1'$ непрерывны: линейной краевой линзы при перпендикулярном торце нет. Во втором порядке непрерывны канонические импульсы, а связь $p_x$ с наклоном внутри и снаружи разная, $p_x=x'-(2/\rho)xx'$ внутри и $p_x=x'$ снаружи. Поэтому $D_2$ непрерывно, а $D_2'$ по разные стороны границы различается: на входе $D_2^{\prime\mathrm{in}}-D_2^{\prime\mathrm{out}}=2D_1D_1'/\rho$, на выходе --- с обратным знаком. Это граничное условие для $D_2$; среднее $\langle D_2/\rho\rangle_C$ входит в $\eta_1^{\mathrm{E}}$ \eqref{eq:eta1-es}. Протяжённое краевое поле электродов здесь не учитывается.

\subsection{Запаздывание за оборот: коэффициенты $\eta_0^{\mathrm{E}}$ и $\eta_1^{\mathrm{E}}$}

Замкнутая орбита
\begin{equation}
\label{eq:x-co}
x=D_1\delta+D_2\delta^2
\end{equation}
определена по внешнему импульсу $\delta$. Период обращения $T=\oint\mathrm{d}l/v$, поэтому
\begin{equation}
\label{eq:T-oint}
1+\tau=\frac{v_s}{C_s}\oint\frac{\mathrm{d}l}{v}
=\Bigl\langle\frac{\mathrm{d}l}{\mathrm{d}s}\,\frac{v_s}{v}\Bigr\rangle_C,
\end{equation}
где
\begin{equation}
\label{eq:avg-C}
\langle f\rangle_C=\frac{1}{C_s}\oint f\,\mathrm{d}s
\end{equation}
--- среднее вдоль референсной орбиты (в отличие от среднего $\langle f\rangle$ по синхротронному периоду). Элемент пути при $y=0$, как и в \eqref{eq:dl},
\begin{equation}
\label{eq:dl-es}
\frac{\mathrm{d}l}{\mathrm{d}s}
=1+\frac{x}{\rho}+\frac{x'^2}{2}+O(x^3).
\end{equation}
На орбите \eqref{eq:x-co} потенциал
\begin{equation}
\label{eq:psi-co}
\chi
=\frac{D_1}{\rho}\,\delta
+\Bigl[\frac{D_2}{\rho}-\frac{n}{2}\Bigl(\frac{D_1}{\rho}\Bigr)^2\Bigr]\delta^2
+O(\delta^3),
\end{equation}
и произведение \eqref{eq:dl-es} на \eqref{eq:vsv-es} равно
\begin{equation}
\label{eq:integrand}
\begin{aligned}
\frac{\mathrm{d}l}{\mathrm{d}s}\frac{v_s}{v}
&=1+\Bigl[
\Bigl(1+\frac{1}{\gamma_s^2}\Bigr)\frac{D_1}{\rho}
-\frac{1}{\gamma_s^2}
\Bigr]\delta
\\
&\quad+\Bigl[
\frac{D_1'^2}{2}
+\Bigl(1+\frac{1}{\gamma_s^2}\Bigr)\frac{D_2}{\rho}
+\frac{5-n}{2\gamma_s^2}\Bigl(\frac{D_1}{\rho}\Bigr)^2
\\
&\qquad\qquad
-\frac{4}{\gamma_s^2}\frac{D_1}{\rho}
+\frac{3\gamma_s^2-1}{2\gamma_s^4}
\Bigr]\delta^2
+O(\delta^3).
\end{aligned}
\end{equation}
Коэффициент $(5-n)/2$ собирается из квадратичного члена потенциала ($-n/2$), из произведения удлинения дуги на линейное замедление ($+1$) и из квадратичного члена ряда скорости ($+3/2$). С геометрическим коэффициентом уплотнения
\begin{equation}
\label{eq:alpha0-es}
\alpha_0=\Bigl\langle\frac{D_1}{\rho}\Bigr\rangle_C,
\end{equation}
вычисленным на дисперсии \eqref{eq:cmp-D-es}, запаздывание имеет вид
\begin{equation}
\label{eq:tau-es-d}
\tau=\eta_0^{\mathrm{E}}\delta+\eta_1^{\mathrm{E}}\delta^2,
\end{equation}
\begin{align}
\eta_0^{\mathrm{E}}
&=\Bigl(1+\frac{1}{\gamma_s^2}\Bigr)\alpha_0-\frac{1}{\gamma_s^2},
\label{eq:eta0-es}\\
\eta_1^{\mathrm{E}}
&=\frac12\langle D_1'^2\rangle_C
+\Bigl(1+\frac{1}{\gamma_s^2}\Bigr)
\Bigl\langle\frac{D_2}{\rho}\Bigr\rangle_C
+\frac{5-n}{2\gamma_s^2}
\Bigl\langle\Bigl(\frac{D_1}{\rho}\Bigr)^2\Bigr\rangle_C
-\frac{4\alpha_0}{\gamma_s^2}
+\frac{3\gamma_s^2-1}{2\gamma_s^4}.
\label{eq:eta1-es}
\end{align}
Индекс $\mathrm{E}$ отмечает коэффициенты электростатического кольца; ниже $\eta_0^{\mathrm{E}}\neq 0$.

\paragraph{Два источника отличий от магнитного кольца.}
Поправки порядка $1/\gamma_s^2$ входят в $\eta^{\mathrm{E}}$ из двух независимых мест. Первое --- орбита: дисперсия удовлетворяет \eqref{eq:cmp-D-es}, а не \eqref{eq:D1}, поэтому $\alpha_0$ отличается от магнитного кольца с теми же квадруполями. Второе --- время пролёта по уже построенной орбите: на смещённом радиусе частица медленнее. В
\begin{equation}
\label{eq:eta0-es-split}
\eta_0^{\mathrm{E}}
=\alpha_0+\frac{\alpha_0}{\gamma_s^2}-\frac{1}{\gamma_s^2}
\end{equation}
слагаемое $\alpha_0$ --- удлинение орбиты (первое место), $\alpha_0/\gamma_s^2$ --- среднее замедление на потенциале (второе место), $-1/\gamma_s^2$ --- зависимость скорости от $\delta$ при нулевом потенциале, общая с магнитным кольцом. В $\eta_1^{\mathrm{E}}$ средние $\frac12\langle D_1'^2\rangle_C$ и $\langle D_2/\rho\rangle_C$ относятся к орбите, а множитель $1/\gamma_s^2$ при $\langle D_2/\rho\rangle_C$, член с $(5-n)$ и член $-4\alpha_0/\gamma_s^2$ --- к времени пролёта.

В нерелятивистском пределе $\eta_0^{\mathrm{E}}\to 2\alpha_0-1$, что согласуется с измерениями на электростатическом кольце CSR~\cite{Grieser2017} (там коэффициент определён через частоту обращения и имеет противоположный знак).

\subsection{Система уравнений и равновесный сдвиг}

При $\sin\varphi\approx\varphi$
\begin{equation}
\label{eq:system-es}
\left\{
\begin{aligned}
\dot\varphi
&=\omega_\rf\bigl(\eta_0^{\mathrm{E}}\delta+\eta_1^{\mathrm{E}}\delta^2\bigr),\\
\dot\delta
&=K\cos\varphi_s\,\varphi\Bigl(1-\frac{\delta}{\gamma_s^2}\Bigr).
\end{aligned}
\right.
\end{equation}
Это система \eqref{eq:system} с заменой $\eta_0\to\eta_0^{\mathrm{E}}$, $\eta_1\to\eta_1^{\mathrm{E}}$. Устойчивая фаза и частота малых колебаний даются \eqref{eq:phis}, \eqref{eq:ws}, а решение --- \eqref{eq:delta-sol} с той же заменой. Условие $\langle\tau\rangle=0$ даёт
\begin{equation}
\label{eq:deq-es}
\delta_\eq=-\frac{\eta_1^{\mathrm{E}}\delta_m^2}{2\eta_0^{\mathrm{E}}}.
\end{equation}

Вне дефлекторов лоренц-фактор связан с $\delta$ рядом \eqref{eq:gamma-delta}, поэтому
\begin{equation}
\label{eq:gamma-eff-es}
\gamma_{\mathrm{eff}}
=\gamma_s+\beta_s^2\gamma_s\,\delta_\eq
+\frac{\beta_s^2\delta_m^2}{4\gamma_s}
+O(\delta_m^4),
\end{equation}
\begin{equation}
\label{eq:dE-es}
\Delta E_\eq
=\beta_s^2E_s\Bigl(
\delta_\eq+\frac{\delta_m^2}{4\gamma_s^2}
\Bigr)
+O(\delta_m^4 E_s).
\end{equation}
Формула \eqref{eq:dE-es} даёт сдвиг средней полной энергии; вне дефлекторов он совпадает со сдвигом кинетической. Лоренц-фактор внутри дефлектора рассмотрен в разделе~\ref{sec:es-gamma-in}.

Условия применимости --- \eqref{eq:valid} с коэффициентами этого кольца и дополнительно
\begin{equation}
\label{eq:valid-es}
\Bigl|\frac{D_1}{\rho}\Bigr|\delta_m\ll 1,
\end{equation}
т.\,е. малость потенциала в дефлекторах. Это условие существенно: например, при $n=2$ и $D_1''=0$ из \eqref{eq:cmp-D-es} следует $D_1/\rho=\gamma_s^2+1$.

\subsection{Кубический член ВЧ-напряжения}

Запаздывание \eqref{eq:tau-es-d} имеет тот же вид, что в магнитном кольце, а резонатор стоит при нулевом потенциале. Поэтому результаты главы~\ref{sec:mag-phase}, в частности \eqref{eq:deq-ph}, переносятся с заменой $\eta_0\to\eta_0^{\mathrm{E}}$, $\eta_1\to\eta_1^{\mathrm{E}}$.
